\section{模型是否学到物理结构：probe、surprise 与 rollout}

Control success 可能来自任务特定捷径，因此课程增加三种 representation diagnostics：probe 检查物理量是否可解码，surprise 检查环境被突然改变时预测误差是否升高，latent geometry 检查状态轨迹是否形成连续结构。它们共同提高解释力，但仍不能替代真实机器人与长 horizon 验证。

\subsection{Latent probing：可读出不等于被显式使用}

slide 49 用 linear probe 和 MLP probe 从冻结 latent 预测 mass、position、velocity 等 ground-truth quantities。若简单线性 probe 已能恢复变量，说明表示中存在较直接的几何编码；若只有 MLP 能恢复，信息可能以非线性方式缠结。probe 结果不保证 predictor 规划时真的依赖该变量。

\lecturefigure{slide-049.jpg}{物理量 latent probing：linear 与 MLP probe 的可解码性}{CS25 V6 official deck, p.~49}

\[
\hat y=Wz+b,
\qquad
R^2=1-\frac{\sum_n(y_n-\hat y_n)^2}{\sum_n(y_n-\bar y)^2}.
\]
其中，$z$ 是冻结 world-model latent，$y$ 是真实物理量，$W,b$ 是 probe 参数，$R^2$ 衡量可解释方差。训练 probe 时不能更新 encoder，否则测试的是新监督任务学习能力，而不是原表示中已有信息。

\begin{warningbox}{Decodable does not mean causal}
一个变量能从 latent 解码，可能只是与真正决策因素相关；要证明模型使用它，需要对 latent、输入或环境做 intervention，并观察预测与动作如何变化。probe 是 representation evidence 的第一层，不是机制终点。
\end{warningbox}

\subsection{World-model violation：用受控突变测 surprise}

slide 50 在 rollout 中突然改变 cube color 或 teleport cube，再观察 prediction error。颜色变化若不影响动力学，理想 state representation 不应过度惊讶；位置瞬移破坏连续 dynamics，应产生明显 error spike。这个设计把``理解物理''改写成可测的 sensitivity profile。

\lecturefigure{slide-050.jpg}{受控环境突变：颜色改变、瞬移与 latent prediction surprise}{CS25 V6 official deck, p.~50}

\[
\operatorname{Surprise}_t
=\left\|P_\phi(E(o_{t-1}),a_{t-1})-E(o_t^{\text{intervened}})\right\|_2^2.
\]
其中，$o_t^{\text{intervened}}$ 是被受控改变的观测，$E$ 是 encoder，$P$ 是 predictor。若 intervention 改变 dynamics-relevant state，surprise 应上升；若只改变被表示忽略的 nuisance factor，surprise 可能较小。

\teachervoice{00:49:08--00:53:13，Lucas 先用 probe 问 latent 是否编码物理量，再用 sudden environmental change 问模型是否对不合理事件感到 ``surprised''。两种测试分别测可读出性与预测一致性。}

\begin{importantbox}{好的 intervention test 必须有对照}
只看 teleportation error spike 不够，还应与 no perturbation、视觉但非动力学变化、不同幅度干预比较。否则模型可能只是对像素分布变化敏感。slide 50 把颜色变化与瞬移并列，正是为了区分 appearance shift 与 state violation。
\end{importantbox}

\subsection{Latent geometry：轨迹是否保持状态邻近关系}

slide 51 将 Push-T 的物理 state grid 与二维投影后的 latent geometry 对照。若相近物理状态在 latent 中也相近，goal distance 才更可能成为有意义的 planning cost。二维投影会丢信息，因此应关注全局连续性、局部折叠和明显撕裂，而不是把某个漂亮形状当成定量证明。

\lecturefigure{slide-051.jpg}{物理状态网格与 learned latent geometry 的相对距离}{CS25 V6 official deck, p.~51}

\begin{knowledgebox}{读图：latent map 支持什么结论}
可支持：表示大致保留某些状态邻近与方向结构。不能支持：每个轴有固定物理语义、所有任务都共享该几何、欧氏距离必然等于控制代价。更可靠的验证应结合 neighborhood preservation、probe、planning success 与 intervention sensitivity。
\end{knowledgebox}

\subsection{Open-loop rollout：长 horizon 会怎样累积误差}

最后一张分析图比较 context input 与 open-loop prediction。前几步往往保持对象与方向，随后误差逐渐累积，尤其在接触、碰撞与多体互动后。这个现象解释了为何当前方法更适合短 horizon MPC：闭环重新观察可以不断重置 latent，纯 open-loop world simulation 则会离开训练分布。

\lecturefigure{slide-052.jpg}{Context-conditioned 与 open-loop latent rollout 的长时行为}{CS25 V6 official deck, p.~52}

\teachervoice{00:53:13--00:55:14，课堂用 latent distance 与 open-loop rollout 分析表示结构，同时没有回避 compounding error。模型在短期保持合理轨迹，不等于已经解决长期 hierarchical planning。}

\subsection{本章小结}

probe、surprise 和 geometry 分别从可解码信息、干预响应和状态空间结构审视 latent。三者与 control 结果互补：一个能规划的模型可能表示难解释，一个 probe 很高的模型也可能不会用这些信息。最可信的判断来自任务成功、机制干预、表征诊断和失败案例的交叉验证。

